% 数据预处理流程；所有缩写在图内释义。
\documentclass[tikz,border=0pt]{standalone}
\usepackage[UTF8,fontset=windows]{ctex}
\usepackage{amsmath}
\setmainfont{Times New Roman}
% 方法流程图的共用样式：黑白矩形框、直角箭头、宋体正文、侧边竖排分区标题。
% 坐标单位为毫米，y 轴向下；版心宽 158 mm，侧边标题占 0--6 mm，正文框占 9--158 mm。
\usetikzlibrary{arrows.meta,calc}
% ctex 默认行距系数为 1.3，图内统一取 1。
\newcommand{\fcsize}[2]{\linespread{1}\fontsize{#1}{#2}\selectfont}
\tikzset{
  fc/.style={x=1mm,y=-1mm,>=Latex,line width=0.45pt},
  box/.style={draw,rectangle,anchor=north west,inner sep=1.5mm,outer sep=0pt,align=left,
    font=\songti\fcsize{8.6}{11}},
  % 固定正文高度并顶端对齐，使同一行各框的标题对齐；框外高为 h + 3 mm。
  h/.style={execute at begin node={\begin{minipage}[t][#1][t]{\linewidth}\raggedright},
    execute at end node={\end{minipage}}},
  io/.style={box,dashed},
  key/.style={box,line width=0.9pt},
  edge/.style={->,line width=0.45pt},
  note/.style={font=\songti\fcsize{8}{9.5},fill=white,inner sep=0.6mm},
}
% 框标题
\newcommand{\bt}[1]{{\heiti #1}\\}
% 侧边分区：\sidelabel{上沿 y}{下沿 y}{竖排文字}
\newcommand{\sidelabel}[3]{%
  \draw (0.3,#1) rectangle (6.1,#2);
  \node[align=center,font=\heiti\fcsize{9.5}{11.5},inner sep=0pt]
    at (3.2,{(#1+#2)/2}) {#3};}
% 总线连接：上方节点列表、下方节点列表、总线 y、总线左端 x、右端 x
\newcommand{\bus}[5]{%
  \foreach \u in {#1} {\draw (\u.south) -- (\u.south |- 0,#3);}
  \draw (#4,#3) -- (#5,#3);
  \foreach \l in {#2} {\draw[edge] (\l.north |- 0,#3) -- (\l.north);}}

\begin{document}
\begin{tikzpicture}[fc]
\path[use as bounding box] (0,0) rectangle (158,155);
\sidelabel{0}{45}{数\\据\\核\\验}
\node[io,text width=146mm,h=16mm] (in) at (9,0) {\bt{输入与目标}%
附件1：250款啤酒、231个化学及理化特征；附件2：相关矩阵；附件3：0--6分记分卡；附件4：50维标准化均值\\
附件8：部分化合物的校准浓度与气味阈值；Schreurs等（2024）数据记录新版：逐品评员0--6分原始评分};
\node[box,text width=69.5mm,h=20mm] (id) at (9,22) {\bt{样本与评分对齐}%
原始评分2197条；\texttt{beer\_name}对应附件1、4的行号\\
按行号对齐后，原始均分与附件4逐维秩相关中位数0.969};
\node[box,text width=69.5mm,h=20mm] (target) at (85.5,22) {\bt{评分尺度判定}%
附件4是按“品评员$\times$维度”标准化再汇总的均值，可复现90.2\%单元\\
标准化丢失0--6分的位置与尺度；主目标为原始品评小组均分};
\bus{in}{id,target}{19}{45.25}{121.75}
\sidelabel{50}{94}{特\\征\\处\\理}
\node[box,text width=44mm,h=20mm] (scale) at (9,50) {\bt{231列的尺度}%
37列经核实为常用对数浓度；33列为线性尺度\\
157列无绝对浓度校准；4列为综合指标};
\node[box,text width=44mm,h=20mm] (na) at (60,50) {\bt{73个化学空值}%
涉及9列、63款；主设定用训练折列均值填充\\
中位数、最低位和原生缺失路由作对照};
\node[box,text width=44mm,h=20mm] (floor) at (111,50) {\bt{未检出状态与极端值}%
142个信号列共享下限，按并列最低秩处理\\
极端值保留；不将下限反推为精确浓度};
\bus{id,target}{scale,na,floor}{47}{32.5}{134.5}
\node[box,text width=146mm,h=13mm] (rule) at (9,78) {\bt{变量用途}%
预测可保留综合指标；归因和生成排除5个综合指标及1个实验内标，留下225个单体化合物};
\bus{scale,na,floor}{rule}{74}{32.5}{134.5}
\sidelabel{100}{155}{三\\问\\接\\口}
\node[io,text width=44mm,h=31mm] (q1) at (9,100) {\bt{问题一：等级预测}%
231列化学特征预测原始品评小组均分，再输出0--6等级\\
验证按25个品评日期组隔离};
\node[io,text width=44mm,h=31mm] (q2) at (60,100) {\bt{问题二：关联归因}%
附件2构建共表达图；225个候选与原始感官均分转为秩\\
对49个维度各报告3个关联候选及稳定性};
\node[io,text width=44mm,h=31mm] (q3) at (111,100) {\bt{问题三：生成与迁移}%
225列映入潜在高斯空间；综合指标另行派生\\
迁移检查9列化学缺陷与6个感官缺陷维度};
\bus{rule}{q1,q2,q3}{97}{32.5}{134.5}
\foreach \u in {q1,q2,q3} {\draw (\u.south) -- (\u.south |- 0,137);}
\draw (32.5,137) -- (134.5,137);
\draw[edge] (83.5,137) -- (83.5,140);
\node[box,text width=146mm,h=12mm] at (9,140) {\bt{跨问题验证原则}%
填充与统计变换只在训练折估计；化学分布按留出啤酒验证，感官预测按日期组验证。};
\end{tikzpicture}
\end{document}
